% ===========================================================================
% Verwendung von KI-Werkzeugen.
% Aufbau übernommen aus docs/jonas_master_thesis.pdf (S. iv).
% @TODO Liste ergänzen, falls weitere Werkzeuge oder Einsatzzwecke hinzukommen.
% ===========================================================================
\clearpage
\thispagestyle{plain}
\setheader{Verwendung von KI-Werkzeugen}

\chapter*{Verwendung von KI-Werkzeugen}

\noindent
Bei der Erstellung dieser Arbeit wurden folgende KI-gestützte Werkzeuge eingesetzt:

\medskip

\begingroup
\rowcolors{1}{white}{white}
\renewcommand{\arraystretch}{1.15}
\noindent
\begin{tabularx}{\textwidth}{@{\hspace{1.5em}}p{0.28\textwidth} X@{}}

Anthropic Claude \newline Opus 4.8 \& Sonnet 4.6 &
Zur Generierung von Quellcode sowie von Analyse- und Visualisierungscode. \newline
Zur Umsetzung der automatisierten Qualitätsbewertung. \newline
Zur Korrektur von \LaTeX-Code. \\[2.2em]

Anthropic Claude \newline Opus 4.8 &
Zur Auf- und Nachbereitung der qualitativen Experteninterviews. \newline
Zur Erstellung einzelner Textpassagen sowie zu deren Rechtschreibkorrektur und
Umformulierung. \newline
Zur Unterstützung bei der Literaturrecherche und bei der Auffindung von
Dokumentationsressourcen zu den verwendeten Standards und Frameworks. \\[2.2em]

Anthropic Claude \newline Fable 5 &
Zur Generierung von Quellcode. \newline
Zur Erstellung einzelner Textpassagen. \newline
Zur Unterstützung bei der Literaturrecherche und bei der Auffindung von
Dokumentationsressourcen zu den verwendeten Standards und Frameworks. \\[2.2em]

OpenAI \newline GPT-5.5 &
Zur Umsetzung der automatisierten Qualitätsbewertung. \newline
Zur Unterstützung bei der Literaturrecherche und bei der Auffindung von
Dokumentationsressourcen zu den verwendeten Standards und Frameworks. \\[2.2em]

DeepL &
Zur Unterstützung bei der Übersetzung von Textpassagen. \\

\end{tabularx}
\endgroup

\bigskip

\noindent
Alle durch KI-Werkzeuge unterstützten Inhalte wurden vom Verfasser kritisch geprüft,
fachlich validiert und bei Bedarf überarbeitet. Die Verantwortung für die endgültige
Fassung liegt vollständig beim Verfasser.

\pagebreak